\xchapter{基于层级语义建模与全局主题约束的免训练视频摘要方法}{Training-Free Video Summarization Method Based on Hierarchical Semantic Modeling and Global Topic Constraints}

已有视频摘要方法虽在公开数据集上取得了一定进展，但仍普遍面临底层视觉特征难以表征高层语义、对大规模人工标注依赖较强以及视频层级结构利用不充分等问题。为应对上述挑战，本章提出一种基于大语言模型的免训练视频摘要方法。该方法通过视频结构化分段、帧级视觉描述生成以及场景级与视频级文本逐层聚合，将视频内容映射至统一的文本语义空间，并利用大语言模型的零样本推理能力完成全局主题提取与场景级语义相关性评估，结合全局主题约束与局部视觉证据实现帧级重要性估计。本章首先分析现有方法的不足，然后概述所提框架的总体结构，并对场景划分、多级描述生成、场景重要性评估以及帧级分数增强等模块进行详细阐述。最后通过对比实验与消融实验验证方法的有效性，并对实验结果进行分析。

\xsection{问题提出}{Problem Statement}

视频摘要任务的目标是在严格长度约束下，从长视频中保留能够表征主要事件与核心叙事的视频子集。视频摘要任务涉及关键帧或关键片段选择、视觉冗余压缩、时序结构建模与高层语义理解等多个方面，是一个复杂的联合优化问题。视频摘要的一般处理流程如图~\ref{fig:video_summary_pipeline}所示，通常包括原始视频输入、帧或场景级预处理、重要性建模、长度约束下的关键内容选择以及最终摘要输出等环节。

\begin{figure}[htbp]
  \centering
  \includegraphics[width=\textwidth]{c3/video_summarization_pipeline_v2.pdf}
  \caption{视频摘要流程示意图}\label{fig:video_summary_pipeline}
\end{figure}

对于开放域视频而言，场景类别、拍摄视角、运动方式和事件结构变化较大，只依赖底层视觉特征往往难以准确把握视频真正的语义重点。针对上述问题，现有视频摘要框架主要有以下不足：1）基于低层视觉特征的无监督方法难以跨越视觉表征与语义表达之间的鸿沟，对于事件是否关键、场景是否承载主题推进等高层问题缺乏刻画及洞察能力。当不同语义事件具有相近视觉外观，或同一事件在不同拍摄条件下呈现显著外观差异时，单纯依赖视觉距离得到的摘要结果容易出现主题偏移与叙事断裂。2）基于监督学习的视频摘要方法对人工标注具有较强依赖。人工标注的视频摘要数据除了成本昂贵外，还具有较强的主观性，同一视频在不同标注者之间往往存在明显偏差，且当测试视频的内容分布与训练数据存在偏移时，模型性能通常显著下降。3）现有方法对视频层级结构的利用仍不充分。视频内容并非简单的帧序列堆叠，而是具有帧（Frame）、场景（Scene）和视频（Video）的层级组织形式。摘要生成不仅需要判断单帧是否清晰或显著，还需要判断该帧所在场景是否承担关键事件表达，该场景与整段视频的主题是否一致。若忽略这种由局部到全局的层级依赖关系，则容易保留视觉上突出但语义无关的片段，或遗漏持续时间较短但叙事价值较高的关键瞬间。

针对上述问题，本文提出一种基于大语言模型的免训练视频摘要方法。该方法不再将摘要生成仅视为低层视觉特征上的代表性筛选问题，而是强调视频内容的层级语义组织与全局主题一致性，通过引入大语言模型的语义理解与推理能力，在无需任何训练的条件下提升传统视频摘要方法难以顾及的语义表达能力与跨场景适应能力。

\xsection{模型结构}{Model Architecture}

本文提出一种基于大语言模型的免训练视频摘要框架，其整体结构如图~\ref{fig:architecture_framework}所示。对于输入视频序列 $\mathcal{V}=\{f_{1},f_{2},\dots,f_{T}\}$，该方法首先依据时间连续性与内容变化将其组织为场景集合 $\mathcal{S}=\{s_{1},s_{2},\dots,s_{K}\}$，随后通过视觉语言模型（Vision-Language Model, VLM）与大语言模型协同构建从帧级描述 $Cap_{f}$ 到场景级描述 $Cap_{s}$ 再到视频级描述 $Cap_{v}$ 的层级语义表示，并在此基础上估计场景相关性分数与帧级重要性分数，最终在摘要长度约束 $L$ 下生成视频摘要子集 $\mathcal{Sum}^{*}\subset \mathcal{V}$。该框架将视频摘要由单纯的视觉代表性筛选转化为面向全局主题一致性的层级语义建模过程，在无需任务特定数据集训练的条件下实现了摘要生成与语义推理的统一。

\begin{figure}
  \centering
  \includegraphics[width=\textwidth]{c3/architecture_v2.pdf}
  \caption{基于大语言模型的免训练视频摘要框架结构图}\label{fig:architecture_framework}
\end{figure}

\xsubsection{视频场景划分模块}{Video Scene Segmentation Module}
在视频摘要任务中，直接在帧级序列上建模会引入大量冗余计算。为此，本文采用核时序分割算法（KTS）\cite{potapov2014category}将输入视频序列 $\mathcal{V}=\{f_{1},f_{2},\dots,f_{T}\}$ 划分为若干时间连续、语义相近的场景单元，作为后续处理的基本单元。核时序分割算法的处理流程如图~\ref{fig:kts}所示。设与视频帧对应的视觉特征序列为 $\mathbf{X}=\{x_{1},x_{2},\ldots,x_{T}\}$（$x_t\in\mathbb{R}^{d}$）。KTS 首先基于核函数 $\kappa(\cdot,\cdot)$ 构建帧间核矩阵 $\mathbf{K}\in\mathbb{R}^{T\times T}$，用于表示视频帧在时间序列中的内容相似关系，随后以段内方差最小化为准则，通过动态规划搜索全局最优边界集合 $\mathcal{B}=\{b_{0},b_{1},\ldots,b_{K}\}$（$0=b_{0}<b_{1}<\cdots<b_{K}=T$），其优化目标为
\begin{equation}
  \min_{\mathcal{B}} J(\mathcal{B})=\sum_{k=1}^{K}\mathrm{Var}(s_k)+\lambda K,
\end{equation}
式中，$\mathrm{Var}(s_k)$——第 $k$ 个场景内部的特征离散程度；$\lambda$——控制分段数量的正则化系数；$K$——场景数量。基于该算法，视频被划分为 $K$ 个连续场景 $\mathcal{S}=\{s_{1},s_{2},\ldots,s_{K}\}$，其中第 $k$ 个场景对应帧区间 $s_k=[b_{k-1}+1,\,b_k]$。相较于基于固定阈值的局部检测方法，KTS 在全局优化意义下兼顾了片段内部一致性与边界显著性。

\begin{figure}[th]
  \centering
  \includegraphics[width=\linewidth]{c3/model_kts.pdf}
  \caption{KTS视频场景划分算法流程}\label{fig:kts}
\end{figure}

通过核时序分割算法，原始长视频被组织为场景集合 $\mathcal{S}$。在此基础上，后续模块能够以场景为基本单元聚合帧级描述并开展语义相关性分析，从而在减少逐帧处理冗余计算的同时保持视频的时序连贯性，为后续场景理解与摘要生成提供结构化输入。


\xsubsection{多级视频语义描述生成模块}{Multi-level Video Semantic Description Generation Module}
在完成场景划分之后，视频已被组织为具有时间连续性的场景单元，但其信息载体仍主要表现为视觉特征或原始帧序列，难以直接用于大语言模型的主题归纳与语义推理。为此，本文设计多级视频语义描述生成模块，将视频内容沿着帧级描述生成、场景级语义聚合和视频级全局概括三层路径逐步映射为结构化文本表示，其整体流程如图~\ref{fig:caption_aggregation}所示。该模块的核心目标是在保留视频时序层次结构的同时，将离散视觉内容转化为具有明确含义的文本单元，从而为后续场景重要性评估提供统一、可比较的文本语义表征。

设经 KTS 划分后的视频包含 $K$ 个场景，场景集合记为 $\mathcal{S}=\{s_1,s_2,\ldots,s_K\}$。对于第 $k$ 个场景 $s_k$，其采样帧集合记为 $\mathcal{F}_k=\{f_{t_{k,1}},f_{t_{k,2}},\ldots,f_{t_{k,n_k}}\}$，其中 $f_{t_{k,i}}$ 表示场景 $s_k$ 中的第 $i$ 个采样帧，$n_k$ 表示该场景的采样帧数量。基于上述层级结构，本文定义帧级描述、场景级描述与视频级描述三类语义变量，分别记为 $Cap_f$、$Cap_s$ 和 $Cap_v$。其中，$Cap_f(t)$ 表示第 $t$ 帧的局部语义描述，$Cap_s(k)$ 表示第 $k$ 个场景的语义概括，$Cap_v$ 表示整段视频的全局主题摘要。

\subsubsection{帧级语义描述生成}
帧级语义描述生成旨在为采样帧构建细粒度的文本语义表征，以显式刻画画面中的主体对象、关键动作与场景环境。对于任意采样帧 $f_t$，记其对应图像输入为 $I_t$，采用视觉语言模型生成自然语言描述，则帧级描述可写为
\begin{equation}
  Cap_f(t)=\mathrm{VLM}(I_t;\theta_{\mathrm{VLM}}),
\end{equation}
式中，$Cap_f(t)$——第 $t$ 帧的语义描述；$I_t$——第 $t$ 帧对应的图像输入；$\mathrm{VLM}(\cdot)$——视觉语言模型；$\theta_{\mathrm{VLM}}$——视觉语言模型参数。通过该映射，原始图像被转换为显式的语言语义单元，使得后续模块不再直接依赖高维视觉特征，而能够在统一的文本语义空间中建模视频内容。需要指出的是，单帧描述虽然能够准确反映局部画面内容，但由于缺乏时间上下文，其对动作延续关系和场景整体语义的表达仍然有限，因此有必要进一步进行上层聚合。

\begin{figure}
  \centering
  \includegraphics[width=\linewidth]{c3/model_caption_aggregation.pdf}
  \caption{多级视频语义描述生成模块流程图}\label{fig:caption_aggregation}
\end{figure}

\subsubsection{场景级语义描述聚合}
场景级语义描述聚合旨在将同一场景内离散分布的多帧语义描述整合为统一的场景级文本表示。对于场景 $s_k$，该过程可分为描述序列构建与场景语义生成两个步骤：首先将采样帧对应的帧级描述按照时间顺序组织为描述序列 $\mathcal{T}_k$，随后将该序列输入场景级语义聚合模型，生成对应的场景级语义描述 $Cap_s(k)$，其形式化表示如下
\begin{equation}
  \mathcal{T}_k=\big[Cap_f(t_{k,1}),Cap_f(t_{k,2}),\ldots,Cap_f(t_{k,n_k})\big].
\end{equation}
\begin{equation}
  Cap_s(k)=\mathrm{LLM}_{\mathrm{scene}}(\mathcal{T}_k;\theta_{\mathrm{LLM}}).
\end{equation}
式中，$\mathcal{T}_k$——第 $k$ 个场景的帧级描述序列；$Cap_f(t_{k,i})$——场景 $s_k$ 内第 $i$ 个采样帧的描述；$n_k$——场景 $s_k$ 的采样帧数量；$Cap_s(k)$——第 $k$ 个场景的语义概括；$\mathrm{LLM}_{\mathrm{scene}}(\cdot)$——场景级语义聚合模型；$\theta_{\mathrm{LLM}}$——大语言模型参数。上述建模方式的关键在于将分散在不同时间位置的帧级描述按时间顺序组织为有序文本序列，使场景内部的语义信息以统一形式输入大语言模型。该序列化表示既保留了场景内部的时间顺序，也为后续语义归纳提供了明确的上下文基础。在具体实现中，本文采用表~\ref{tab:scene_prompt_example} 所示的场景级语义聚合提示词模板，将同一场景内的帧级描述结合系统提示词与用户指令，引导大语言模型按时间顺序生成对应的场景级语义描述。该提示词模板进一步约束模型的关注重点，使其围绕场景中的主要事件、参与主体及变化过程进行归纳，从而减少偏离场景核心内容的冗余输出。

\begin{table}[htbp]
  \centering
  \caption{场景级语义聚合提示词示例}\label{tab:scene_prompt_example}
  \begin{tblr}{
    width=\linewidth,
    colspec={X[1.1,c]X[4.9,l]},
    rowsep=2pt,
    cells={valign=m},
  }
    \toprule
    \textbf{角色} & \textbf{提示词内容}\\
    \midrule
    系统 & You are a video scene summarization assistant. Given a sequence of frame captions from the same scene, generate one concise scene-level summary that captures the main event, the key subjects, and the temporal progression of the scene. \\
    \midrule
    用户 & Frame captions: (1) A man opens a car door and steps out with a suitcase. (2) He walks toward the entrance of a hotel. (3) He pauses briefly to check his phone before entering the lobby. Task: Summarize this scene in one sentence. \\
    \midrule
    模型输出 & A man arrives at a hotel with a suitcase, walks to the entrance, and enters the lobby after briefly checking his phone. \\
    \bottomrule
  \end{tblr}
\end{table}

与单帧描述相比，$Cap_s(k)$ 不再仅反映某一时刻的局部内容，而是对整个场景中的连续语义信息进行综合表示，不同帧之间的互补信息在聚合过程中得到整合，场景级描述也能够更稳定地反映场景中的核心事件、主要参与主体、动作变化及整体语义；同时，场景内部的冗余细节和偶然噪声也得到抑制，从而形成更完整、更连贯的场景语义表示，并为后续不同场景之间的重要性比较提供更稳定的基础。进一步地，这种由多帧描述生成的场景级表示受单帧采样偏差和瞬时视觉变化的影响较小，因此在不同场景之间具有更好的可比性。基于这一表示，后续模块能够在统一的语义层面开展场景重要性评估，从而使视频级主题概括与场景级评分之间保持更稳定的对应关系。

\subsubsection{视频级全局语义摘要}
在获得全部场景级描述后，需要进一步从全局视角对整段视频的主题内容进行归纳，以形成后续场景评价的统一语义参照。设视频包含 $K$ 个场景，其场景级描述序列记为 $\{Cap_s(1),Cap_s(2),\ldots,Cap_s(K)\}$，则视频级语义摘要可表示为
\begin{equation}
  Cap_v=\mathrm{LLM}_{\mathrm{video}}\big(\{Cap_s(1),Cap_s(2),\ldots,Cap_s(K)\};\theta_{\mathrm{LLM}}\big).
\end{equation}
式中，$Cap_v$——整段视频的全局语义摘要；$\mathrm{LLM}_{\mathrm{video}}(\cdot)$——视频级摘要模型；$Cap_s(k)$——第 $k$ 个场景的语义描述；$K$——场景总数。
上述过程的关键在于将按时间顺序排列的场景级描述作为统一输入，使大语言模型能够在跨场景范围内综合不同片段的事件信息、主题线索与叙事关系，生成面向整段视频的全局语义摘要。在具体实现中，本文采用表~\ref{tab:video_prompt_example} 所示的视频级摘要提示词模板，将全部场景级描述输入大语言模型，并要求模型围绕视频主题、核心事件与整体叙事线索进行归纳，输出对应的视频级语义摘要。

\begin{table}[htbp]
  \centering
  \caption{视频级全局语义摘要提示词示例}\label{tab:video_prompt_example}
  \begin{tblr}{
    width=\linewidth,
    colspec={X[1.1,c]X[4.9,l]},
    rowsep=2pt,
    cells={valign=m},
  }
    \toprule
    \textbf{角色} & \textbf{提示词内容}\\
    \midrule
    系统 & You are a video summarization assistant. Given a sequence of scene-level summaries from the same video, generate one concise video-level summary that captures the main theme, the key events, and the overall narrative flow. \\
    \midrule
    用户 & Scene summaries: (1) A man arrives at a hotel and checks in. (2) He visits several landmarks in the city center. (3) He watches the sunset from a riverside viewpoint and returns to the hotel. Task: Summarize the whole video in one or two sentences. \\
    \midrule
    模型输出 & The video presents a traveler's city trip, covering hotel arrival, landmark sightseeing, and an evening riverside sunset, and forms a complete urban travel narrative. \\
    \bottomrule
  \end{tblr}
\end{table}

与前两级表示不同，视频级描述面向整段视频的主题归纳，其关注点不再局限于单一场景的局部内容，而是对视频中的核心事件、叙事主线及整体语义进行集中概括。由于不同场景之间的语义信息在该阶段得到统一整合，$Cap_v$ 能够更稳定地反映视频的总体主题，并减少局部场景差异对后续判断的影响。基于这一全局语义表示，后续模块可以在统一主题条件下评估各场景的重要性，从而使场景级评分与视频整体叙事保持更高的一致性。

综上，本文通过帧级描述生成、场景级语义聚合与视频级全局摘要三个步骤，构建了由局部画面内容到全局主题概括的层级语义表示链路。该模块不仅保持了视频内容由局部到整体的时序组织形式，而且将原始视觉信息转化为适于大语言模型处理的文本表达，为后续基于语义相关性的评分机制奠定了表示基础。

\xsubsection{全局主题约束下的场景语义评分模块}{Global Theme-Constrained Scene Semantic Scoring Module}
在获得场景级描述 $Cap_s(k)$ 与视频级描述 $Cap_v$ 之后，摘要生成的关键问题可进一步表述为：如何在统一语义参照下度量各场景对整段视频主题表达的贡献程度。为此，本文构建基于大语言模型的场景语义评分模块，利用视频全局语义摘要作为评分条件来衡量各场景的重要性，避免仅依赖视觉统计量而导致评分偏离视频主题。

\begin{figure}
  \centering
  \includegraphics[width=0.95\linewidth]{c3/model_llm_query.pdf}
  \caption{全局主题约束下的场景语义评分模块}\label{fig:llm_query}
\end{figure}

设输入视频经 KTS 划分后得到 $K$ 个场景，记场景集合为 $\mathcal{S}=\{s_1,s_2,\ldots,s_K\}$。其中，$Cap_s(k)$ 表示第 $k$ 个场景 $s_k$ 的场景级语义描述，$Cap_v$ 表示整段视频的全局语义摘要，$\theta_{\mathrm{LLM}}$ 表示评分模型参数。本文将第 $k$ 个场景的重要性表示为：在给定全局主题条件下，该场景被纳入最终摘要的条件概率，可表示为
\begin{equation}
  S_{\mathrm{LLM}}(k)=p\big(y_k=1\mid Cap_s(k),Cap_v;\theta_{\mathrm{LLM}}\big),
\end{equation}
式中，$S_{\mathrm{LLM}}(k)$——第 $k$ 个场景的重要性分数；$y_k$——场景 $s_k$ 是否应被纳入最终摘要的二值变量；$\theta_{\mathrm{LLM}}$——评分模型参数。
其中，$y_k\in\{0,1\}$ 表示场景 $s_k$ 是否应被纳入最终摘要，$S_{\mathrm{LLM}}(k)\in[0,1]$ 为对应的归一化重要性分数。上述定义表明，该模块并不直接对视觉特征进行度量，而是利用场景语义与视频主题之间的条件一致性完成重要性判断。

为将上述判别问题转化为实际可执行的评分流程，本文对每个场景分别构造一条独立的评分查询。具体而言，对于第 $k$ 个场景，将视频级语义摘要 $Cap_v$ 与对应的场景级描述 $Cap_s(k)$ 按照统一模板组织为提示输入，得到
\begin{equation}
  q_k=\Pi\big(Cap_v,Cap_s(k)\big), \qquad k=1,2,\ldots,K,
\end{equation}
以将视频全局主题、当前场景内容以及评分任务要求整合为统一的文本输入。对全部场景重复该构造后，可得到与场景集合一一对应的提示序列 $\{q_1,q_2,\ldots,q_K\}$。为便于说明具体格式，表~\ref{tab:scene_score_prompt_example} 给出了单个场景的提示词示例。

\begin{table}[htbp]
  \centering
  \caption{场景语义评分提示词示例}\label{tab:scene_score_prompt_example}
  \begin{tblr}{
    width=\linewidth,
    colspec={X[1.1,c]X[4.9,l]},
    rowsep=2pt,
    cells={valign=m},
  }
    \toprule
    \textbf{角色} & \textbf{提示词内容}\\
    \midrule
    系统 & You are a video summarization evaluator. Given the global video summary and the current scene description, output only one importance score from 1 to 100. \\
    \midrule
    用户 & \textbf{Global video summary} = A travel video records the complete visiting process of the Eiffel Tower, including arrival, sightseeing, and panoramic city views. \textbf{Current scene description} = Several tourists are waiting in a long queue near the entrance, with limited change in actions and no key landmark view. Task: Evaluate how important the current scene is for preserving the main theme and narrative completeness of the video. Output score only. \\
    \midrule
    模型输出 & 22 \\
    \bottomrule
  \end{tblr}
\end{table}

在实际执行时，将每个 $q_k$ 单独输入大语言模型，并要求模型仅输出一个介于 0 到 100 之间的整数分值，以表达当前场景对保持视频主题与叙事完整性的贡献程度。记该原始输出为 $r_k$，则其推理过程可写为 $r_k=\mathrm{LLM}(q_k;\theta_{\mathrm{LLM}})$。考虑到模型输出中可能包含额外文本或超出预设范围的数值，本文在后处理阶段首先通过正则匹配提取其中的数值结果，并将其截断至区间 $[0,100]$；随后再除以 100 映射到 $[0,1]$ 区间，得到最终的场景语义分数。基于上述过程，单个场景的评分映射可形式化为
\begin{equation}
  S_{\mathrm{LLM}}(k)=\mathrm{Norm}\big(\mathrm{LLM}(\Pi(Cap_v,Cap_s(k));\theta_{\mathrm{LLM}})\big),
\end{equation}

对全部 $K$ 个场景独立重复这三个步骤后，即可按照场景原始顺序得到视频的场景级评分序列
\begin{equation}
  \mathrm{Score}_{\mathrm{scene}}=\big[S_{\mathrm{LLM}}(1),S_{\mathrm{LLM}}(2),\ldots,S_{\mathrm{LLM}}(K)\big].
\end{equation}
式中，$\mathrm{Score}_{\mathrm{scene}}$——视频的场景级评分序列；$S_{\mathrm{LLM}}(k)$——第 $k$ 个场景的语义重要性分数；$K$——场景总数。

与基于视觉显著性、聚类中心距离或重构误差的传统评分方式不同，上述模块以视频全局主题为条件对场景语义进行判别，因而能够在开放域视频中有效区分视觉上显著但与主题无关的片段，以及视觉变化有限但对叙事不可或缺的片段。然而，由于同一场景内所有帧共享相同的评分值，$\mathrm{Score}_{\mathrm{scene}}$ 本质上仅提供粗粒度的语义先验，尚无法刻画场景内部不同帧之间的质量差异与局部语义贡献。为弥补这一不足，下一小节将在场景级分数的基础上，引入多粒度语义相似度实现帧级重要性增强。

\xsubsection{场景先验引导的帧级语义分数增强模块}{Scene-prior-guided Frame-level Semantic Score Refinement Module}

上一小节得到的场景级分数序列虽然能够从全局主题层面对各场景的重要性作出判断，但其在同一场景内部取值恒定，无法区分帧间的细粒度差异。当场景中同时存在高价值和低价值帧时，仅依赖场景级分数易导致关键帧遗漏或冗余帧误留。为了解决这个问题，本文设计场景先验引导的帧级语义分数增强模块，在继承场景级主题约束的基础上，融合帧与场景、帧与视频两个层面的语义相似度，对各采样帧的重要性进行精细化建模。如图~\ref{fig:score_refinement} 所示，该模块以场景级分数为先验，并结合每一帧与所属场景及整段视频的语义相关性，将场景内部恒定的粗粒度评分细化为逐帧变化的重要性分数。

\begin{figure}
  \centering
  \includegraphics[width=0.95\linewidth]{c3/model_score_refinement.pdf}
  \caption{场景先验引导的帧级语义分数增强模块}\label{fig:score_refinement}
\end{figure}

设视频经采样后得到帧索引集合 $\mathcal{P}=\{p_1,p_2,\ldots,p_N\}$，其中 $p_j$ 为第 $j$ 个采样位置对应的原始帧索引，$N$ 为采样帧总数。进一步定义场景归属映射 $\sigma(j)$，使得 $p_j\in s_{\sigma(j)}$，即 $\sigma(j)$ 返回采样帧 $p_j$ 所在的场景编号。在已有场景语义分数的基础上，本文进一步为每个采样帧构造综合分数 $F(p_j)$，使其同时反映全局主题一致性、场景局部代表性与帧级语义贡献三方面信息。

\subsubsection{统一文本语义嵌入}
为确保帧级描述、场景级描述与视频级描述能够在同一语义空间中进行相似性比较，本文采用共享文本编码器 $\mathrm{Enc}(\cdot)$ 对三类描述统一进行向量化。该编码器由预训练句向量模型提供，输出维度记为 $d$。对于采样帧 $p_j$、其所属场景 $s_{\sigma(j)}$ 以及整段视频，分别定义如下语义嵌入：
\begin{equation}
  \mathbf{e}_{f}(j)=\mathrm{Enc}\big(Cap_f(p_j)\big)\in\mathbb{R}^{d},
\end{equation}
\begin{equation}
  \mathbf{e}_{s}\big(\sigma(j)\big)=\mathrm{Enc}\big(Cap_s(\sigma(j))\big)\in\mathbb{R}^{d},
\end{equation}
\begin{equation}
  \mathbf{e}_{v}=\mathrm{Enc}(Cap_v)\in\mathbb{R}^{d}.
\end{equation}
式中，$\mathbf{e}_f(j)$、$\mathbf{e}_s(\sigma(j))$、$\mathbf{e}_v$——分别表示采样帧 $p_j$、其所属场景 $s_{\sigma(j)}$ 以及整段视频的文本语义嵌入；$Cap_f(p_j)$、$Cap_s(\sigma(j))$、$Cap_v$——分别表示采样帧描述、所属场景描述和视频级全局语义摘要；$\mathrm{Enc}(\cdot)$——共享文本编码器；$d$——嵌入维度；$\sigma(j)$——采样帧 $p_j$ 所在的场景编号。为简化后文记号，令 $\mathbf{e}_{s,j}=\mathbf{e}_{s}(\sigma(j))$。在相似度计算之前，编码器输出向量统一进行 $L_2$ 归一化，所得归一化向量分别记为 $\mathbf{z}_{f,j}$、$\mathbf{z}_{s,j}$ 与 $\mathbf{z}_{v}$，满足 $\|\mathbf{z}_{f,j}\|_2=\|\mathbf{z}_{s,j}\|_2=\|\mathbf{z}_{v}\|_2=1$。共享编码器的作用在于把帧级、场景级和视频级文本描述映射到同一个语义空间中，使不同层级的表示能够直接比较，从而为后续语义相似度计算提供统一基础。

\subsubsection{局部与全局语义相关性建模}
在统一嵌入表示的基础上，本文从局部场景一致性与全局主题一致性两个维度分别度量帧级语义相关性。首先，定义帧与所属场景之间的局部相关性分数为
\begin{equation}
  S_{FS}(p_j)=\mathbf{z}_{f,j}^{\top}\mathbf{z}_{s,j},
\end{equation}
该分数衡量采样帧 $p_j$ 对所属场景语义中心的代表程度：分数越高，表明该帧越能准确概括场景的核心事件与视觉内容。进一步，定义帧与视频全局摘要之间的全局相关性分数为
\begin{equation}
  S_{FV}(p_j)=\mathbf{z}_{f,j}^{\top}\mathbf{z}_{v}.
\end{equation}
与 $S_{FS}(p_j)$ 的局部视角不同，$S_{FV}(p_j)$ 直接从视频叙事主线出发衡量帧的全局契合度，因而能够为持续时间短暂但语义密度较高的关键瞬间提供额外的重要性补偿。

\subsubsection{帧级分数融合与实现流程}
在帧级语义增强阶段，场景分数 $S_{\mathrm{LLM}}(\sigma(j))$ 充当粗粒度先验，约束当前帧所属场景的整体重要性水平；$S_{FS}(p_j)$ 刻画帧对所在场景的代表能力；$S_{FV}(p_j)$ 提供帧与全局主题之间的直接关联证据。综合上述三项信号，本文将采样帧 $p_j$ 的最终重要性分数定义为
\begin{equation}\label{eq:frame_refinement_score}
  F(p_j)=\alpha\,S_{\mathrm{LLM}}\big(\sigma(j)\big)\,S_{FS}(p_j)+S_{FV}(p_j),
\end{equation}
式中，$F(p_j)$——采样帧 $p_j$ 的最终重要性分数；$\alpha$——场景先验调节系数；$S_{\mathrm{LLM}}(\sigma(j))$——该帧所属场景的语义重要性分数；$S_{FS}(p_j)$——帧与场景的局部相关性分数；$S_{FV}(p_j)$——帧与视频主题的全局相关性分数。
其中 $\alpha\geq 0$ 为场景先验调节系数，用于控制场景级语义先验与全局主题匹配项之间的相对贡献。

具体而言，乘积项 $S_{\mathrm{LLM}}(\sigma(j))S_{FS}(p_j)$ 构成受场景先验约束的局部贡献项，既能将场景级重要性细化到帧级分辨率，又能抑制场景内部与主要语义偏离的低质量帧；加性项 $S_{FV}(p_j)$ 则从全局主题层面提供额外补偿，使持续时间较短但与视频主旨高度一致的关键帧仍能获得较高得分。通过这种局部细化与全局补偿的融合方式，最终评分在保持视频叙事主线一致性的同时，提高了对场景内部关键内容的区分能力。基于上述定义，整段视频的帧级评分序列可表示为
\begin{equation}
  \mathbf{f}=\big[F(p_1),F(p_2),\ldots,F(p_N)\big].
\end{equation}

在具体实现中，该模块按三个步骤完成帧级评分：首先根据场景归属映射将场景级分数 $S_{\mathrm{LLM}}(\sigma(j))$ 赋予对应的采样帧，为每帧注入场景层面的主题先验；其次利用文本嵌入逐帧计算帧-场景相似度 $S_{FS}(p_j)$ 与帧-视频相似度 $S_{FV}(p_j)$；最后按式~\eqref{eq:frame_refinement_score} 融合得到最终评分。由于所有输入均来自前序模块且计算过程仅涉及嵌入与余弦相似度，该模块无需额外训练即可直接部署。与场景级评分仅区分不同场景间的重要性差异相比，帧级增强使同一场景内各帧的得分得以拉开差距，从而在最终的摘要阶段为关键帧的精准筛选提供更细粒度的排序依据。

\xsection{实验与分析}{Experiments and Analysis}
\xsubsection{实验细节}{Experimental Details}
\subsubsection{预处理与场景划分}
本文采用 KTS 算法对视频进行场景划分，并将得到的连续片段统一作为后续语义建模的场景单元。为控制计算开销，首先对原始视频进行时间下采样，实验中默认以每秒 2 帧的频率提取候选帧；随后使用 GoogLeNet 提取帧级视觉特征，并据此构建 KTS 所需的时序核矩阵。正则化参数 $\lambda$ 随视频时长自适应调整，以保证划分结果既能保留关键语义边界，又不过度切碎视频结构。实际统计中，所得场景平均时长控制在 3--5 秒之间，能够较好兼顾语义完整性与后续推理成本。

\subsubsection{模型与推理配置}
帧级语义描述由视觉语言模型生成。综合生成质量与推理效率，本文采用 LLaVA-Video-7B-Qwen2 作为视觉描述模型，对每秒 2 帧采样得到的全部采样帧逐帧生成文本描述。在场景级描述聚合与重要性评分阶段，使用 DeepSeek-V3.2 对多帧描述进行归纳；相关提示词均采用 Zero-shot 形式构造，温度参数设置为 0.2，以在稳定性与表达多样性之间取得平衡。对于帧级分数增强模块，本文使用 Qwen3-Embedding-8B 将帧级、场景级与视频级文本描述映射到统一语义空间，并在归一化后计算帧与场景、帧与全局视频主题之间的相似度。

\subsubsection{摘要生成与实验环境}
在得到帧级重要性分数后，本文进一步将其映射回原始时间轴，并结合 KTS 给出的场景边界构造候选摘要片段。对于每个场景片段，以片段持续时长作为约束代价，以片段内部帧级重要性分数的综合统计量作为优化收益，在摘要总时长不超过原视频 15\% 的约束下求解最优片段集合，最终按时间顺序拼接形成二值视频摘要。该优化过程采用 Google OR-Tools 提供的 0/1 背包求解器实现，与 SumMe 和 TVSum 的主流评测协议保持一致，从而保证与现有方法比较时的公平性。

全部实验均在 Ubuntu 22.04 环境下完成，硬件配置为单张 NVIDIA RTX5090 GPU。除基于 Google OR-Tools 实现的摘要选择模块外，其余预处理、特征提取、文本编码与相似度计算流程均基于 PyTorch 实现。由于本文方法属于免训练推理框架，实验中不涉及额外参数训练或任务特定微调，全部结果均在统一硬件环境与一致推理配置下获得。

\xsubsection{对比实验}{Comparative Experiments}
为验证本文方法在公开视频摘要任务中的有效性，本节从定量分析与定性分析两个层面展开对比。前者关注模型在标准数据集上的整体性能表现，后者则进一步考察所生成摘要在语义完整性、冗余抑制与叙事连贯性方面的主观质量。

\subsubsection{$F_1$ 分数定量分析}
表\ref{tab:vs_benchmark} 给出了本文方法与代表性监督方法、无监督方法及免训练方法在 SumMe 和 TVSum 数据集上的对比结果。其中，SumMe 报告 Max $F_1$-Score，TVSum 报告 Avg $F_1$-Score，均与各数据集的主流评估协议保持一致。为便于理解表中各基线方法的性能差异，下面对主要对比方法作简要介绍。

dppLSTM\cite{zhang2016video}：该方法将视频摘要建模为结构化预测问题，采用长短期记忆网络建模视频帧之间的变范围时序依赖关系，以此预测关键帧或关键子镜头，并结合辅助数据集与领域自适应策略缓解标注数据不足带来的训练困难。

ACGAN\cite{he2019unsupervised}：该方法构建了基于注意力条件生成对抗网络的视频摘要框架，由生成器预测帧级重要性分数并生成加权特征，由判别器判别加权特征与原始特征的差异，同时引入条件特征选择器和多头自注意力机制，以增强对长程时序依赖与关键时间区域的建模能力。

MPFN\cite{khan2024deep}：该方法提出了 ViT 辅助的深度金字塔细化网络，首先利用带有视觉 Transformer 主干的密集预测模块提取帧表示，再通过多尺度渐进式特征融合与定制金字塔细化块对多层特征进行逐级增强，最终预测帧级重要性分数并生成视频摘要。

CM-AVS\cite{huang2026conditional}：该方法从条件建模视角出发，引入多个具有语义含义的随机变量及其联合分布来刻画视频摘要中的关键因素，并通过辅助分布改进训练过程，同时设计条件注意力模块以缓解多模态输入带来的性能退化问题。

DR-DSN\cite{zhou2018deep}：该方法将视频摘要视为序列决策过程，利用深度摘要网络为每一帧预测被选中的概率，并基于强化学习进行端到端训练；其奖励函数同时约束摘要结果的多样性与代表性，因此无需人工标注即可完成无监督学习。

OM-AE\cite{rochan2019video}：该方法针对配对标注数据稀缺的问题，提出从非配对数据学习视频摘要的框架，通过对抗目标学习从原视频到摘要视频的映射，并额外引入多样性约束，使生成的摘要结果既符合目标摘要分布，又具备较好的视觉多样性。

SUM-SR\cite{li2024unsupervised}：该方法借鉴生成对抗学习思想但去除了判别器，采用更简单的损失函数，并通过重构器与帧选择器交替迭代训练的策略优化摘要生成过程；此外还在训练和测试阶段引入可训练掩码向量，以提升模型稳定性和训练效率。

SegSum\cite{bang2025integrate}：该方法显式利用 KTS 划分得到的片段信息，将视频片段级时空关系纳入摘要器的学习过程中，并结合集中注意力机制提升对片段内部与片段之间结构信息的建模能力，从而改善传统逐帧打分再汇聚策略的信息损失问题。

M3Sum\cite{wang2024m3sum}：该方法提出了一种模块化多模态免训练视频摘要框架，通过将音频、视频帧等不同模态统一转换为文本描述，并设计无参数的文本对齐机制进行跨模态融合，从而利用大语言模型的长上下文理解能力在无需训练数据和参数更新的条件下生成摘要结果。

\begin{table}[htbp]
  \centering
  \caption{本文方法与代表性视频摘要方法的性能对比}\label{tab:vs_benchmark}
  \begin{tblr}{
    width=\linewidth,
    colspec={X[1.2,l]X[0.9,c]X[0.9,c]X[0.9,c]},
    rowsep=2pt,
  }
    \toprule
    \textbf{方法} & \textbf{类别} & \textbf{SumMe} & \textbf{TVSum}\\
    \midrule
    dppLSTM (2016)\cite{zhang2016video} & 监督 & 38.6 & 54.7 \\
    ACGAN (2019)\cite{he2019unsupervised} & 监督 & 47.2 & 59.4 \\
    MPFN (2024)\cite{khan2024deep} & 监督 & 51.9 & 62.4 \\
    CM-AVS (2026)\cite{huang2026conditional} & 监督 & 52.4 & \textbf{67.5} \\
    \midrule
    DR-DSN (2018)\cite{zhou2018deep} & 无监督 & 41.4 & 57.6 \\
    OM-AE (2019)\cite{rochan2019video} & 无监督 & 37.7 & 51.5 \\
    SUM-SR (2024)\cite{li2024unsupervised} & 无监督 & 51.3 & 60.2 \\
    SegSum (2025)\cite{bang2025integrate} & 无监督 & \textbf{54.0} & 62.0 \\
    \midrule
    M3Sum-SP (2024)\cite{wang2024m3sum} & 免训练 & 43.6 & 56.9 \\
    M3Sum-PCoT (2024)\cite{wang2024m3sum} & 免训练 & 41.9 & 57.6 \\
    \textbf{本章方法} & \textbf{免训练} & \textbf{48.1} & \textbf{61.1} \\
    \bottomrule
  \end{tblr}
\end{table}

\textbf{与免训练方法的比较。}\enspace 在同属免训练范式的方法中，本文方法在 SumMe 上分别高于 M3Sum-SP（43.6\%）和 M3Sum-PCoT（41.9\%）4.5 与 6.2 个百分点，在 TVSum 上分别高于 56.9\% 和 57.6\% 4.2 与 3.5 个百分点。M3Sum 主要通过多模态文本化与无参数对齐机制完成摘要生成，而本文进一步引入场景级主题评分与帧级语义增强，使全局主题约束能够显式传递到帧级选择过程，因此在两个数据集上均表现出更稳定的优势。

\textbf{与无监督方法的比较。}\enspace 本文方法在完全免训练的条件下，于 SumMe 上（48.1\%）大幅超越早期无监督方法 DR-DSN（41.4\%）与 OM-AE（37.7\%），分别取得 6.7 与 10.4 个百分点的提升；在 TVSum 上（61.1\%）也分别领先 DR-DSN（57.6\%）与 OM-AE（51.5\%）3.5 与 9.6 个百分点。上述结果表明，基于视觉重构或自编码器的传统无监督方法在语义建模能力上存在明显瓶颈，而本文通过引入大语言模型的层级语义推理，即使在不进行任何参数学习的前提下，仍能有效弥合视觉特征与高层叙事语义之间的鸿沟。与近年较强的无监督基线相比，SUM-SR（51.3\%/60.2\%）与 SegSum（54.0\%/62.0\%）在 SumMe 上分别高出本文方法 3.2 与 5.9 个百分点，但这两种方法均需在目标数据集上执行完整的训练流程以学习帧间时序依赖与选择策略。值得注意的是，在 TVSum 上本文方法（61.1\%）与 SUM-SR（60.2\%）仅差 0.9 个百分点，表明层级语义评分在主题结构明确的数据集上已能取得与训练式无监督方法相近的帧级区分能力。

\textbf{与监督方法的比较。}\enspace 本文方法在两个数据集上均超越了早期监督基线 dppLSTM（38.6\%/54.7\%），在 SumMe 上还超越了 ACGAN（47.2\%），领先幅度分别为 9.5/6.4 与 0.9 个百分点。与近年最优监督方法 MPFN（51.9\%/62.4\%）和 CM-AVS（52.4\%/67.5\%）相比，本文方法在 SumMe 上分别落后 3.8 与 4.3 个百分点，在 TVSum 上分别落后 1.3 与 6.4 个百分点。然而需要指出的是，MPFN 与 CM-AVS 均依赖于人工标注的帧级或片段级训练信号，并通过端到端参数优化学习帧间时序建模与选择策略，而本文方法不引入任何可学习参数。在这一严格约束下仍能与监督方法保持可比的性能水平，充分印证了将大语言模型的零样本语义理解能力应用于视频摘要任务的可行性与发展潜力。

\subsubsection{秩相关性系数定量分析}

为进一步评价模型输出帧级重要性分数与人工标注之间的排序一致性，表~\ref{tab:rank_correlation_benchmark} 给出了本文方法与代表性视频摘要方法在 SumMe 和 TVSum 数据集上的 Spearman 系数 $\rho$ 与 Kendall 系数 $\tau$ 对比结果。为便于理解表中不同方法在帧级排序建模上的差异，下面对主要对比方法作简要介绍。

VASNet\cite{fajtl2018summarizing}：该方法提出基于软自注意力机制的监督式关键镜头摘要框架，以全连接网络替代双向循环网络完成序列到序列映射，从而在降低模型复杂度的同时建模帧间全局依赖关系。

DSNet\cite{zhu2020dsnet}：该方法提出 Detect-to-Summarize 框架，其中锚框式变体 DSNet-AB 通过多尺度时间兴趣提议实现关键片段定位与评分，锚点自由变体 DSNet-AF 则直接预测帧级重要性与片段位置，以显式引入时间一致性与摘要完整性约束。

DMASum\cite{wang2020query}：该方法围绕 Softmax Bottleneck 问题构建双混合注意力模型，通过二次自查询注意力增强模型表达能力，并结合单帧元学习策略提升小规模数据条件下的泛化性能。

PGL-SUM\cite{apostolidis2021combining}：该方法联合使用全局与局部多头注意力机制，并引入绝对位置编码，以从不同粒度上建模长程帧间依赖关系，从而提升视频帧重要性排序的时序感知能力。

MSVA\cite{ghauri2021supervised}：该方法融合静态视觉特征、运动特征等多种表示，并在特征融合前引入并行注意力机制，以增强对视觉内容与动态信息的联合建模能力。

iPTNet\cite{jiang2022joint}：该方法通过跨任务样本迁移，将视频时刻定位任务中的丰富标注知识迁移到视频摘要任务中，并利用重要性传播模块与协同教学机制改善帧级重要性估计。

CSTA\cite{son2024csta}：该方法将单视频帧特征堆叠为类图像表示，并采用二维卷积建模帧间与帧内关系，以较低计算开销提取时空显著性并预测帧级重要性分数。

LLMVS\cite{lee2025video}：该方法首先利用多模态大语言模型将视频帧转换为文本描述，再基于局部上下文由大语言模型评估帧重要性，并通过全局注意力机制对局部分数进行训练微调统一修正，以兼顾局部细节与整体叙事线索。

\begin{table}[htbp]
  \centering
  \caption{本文方法与代表性视频摘要方法的秩相关性系数对比}\label{tab:rank_correlation_benchmark}
  \begin{tblr}{
    width=\linewidth,
    colspec={X[3.2,l]X[c]X[c]X[c]X[c]},
    rowsep=2pt,
    cell{1}{2}={c=2}{c},
    cell{1}{4}={c=2}{c},
    % hline{3}={2-3}{0.5pt},
    % hline{3}={4-5}{0.5pt},
  }
    \toprule
    \textbf{方法} & \textbf{SumMe} & & \textbf{TVSum} & \\
    \midrule
    & $\rho$ & $\tau$ & $\rho$ & $\tau$ \\
    \midrule
    dppLSTM (2016)\cite{zhang2016video} & 0.049 & 0.040 & 0.055 & 0.042 \\
    VASNet (2018)\cite{fajtl2018summarizing} & 0.170 & 0.160 & 0.170 & 0.160 \\
    DSNet-AB (2020)\cite{zhu2020dsnet} & 0.059 & 0.051 & 0.129 & 0.108 \\
    DSNet-AF (2020)\cite{zhu2020dsnet} & 0.046 & 0.037 & 0.138 & 0.113 \\
    DMASum (2020)\cite{wang2020query} & 0.089 & 0.063 & 0.267 & 0.203 \\
    PGL-SUM (2021)\cite{apostolidis2021combining} & -- & -- & 0.157 & 0.206 \\
    MSVA (2021)\cite{ghauri2021supervised} & 0.230 & 0.200 & 0.210 & 0.190 \\
    iPTNet (2022)\cite{jiang2022joint} & 0.119 & 0.101 & 0.163 & 0.134 \\
    CSTA (2024)\cite{son2024csta} & 0.274 & 0.246 & 0.255 & 0.194 \\
    LLMVS (2025)\cite{lee2025video} & \textbf{0.282} & \textbf{0.253} & 0.275 & \textbf{0.211} \\
    \midrule
    \textbf{本章方法} & 0.267 & 0.218 & \textbf{0.280} & 0.191 \\
    \bottomrule
  \end{tblr}
\end{table}

从表~\ref{tab:rank_correlation_benchmark} 的整体结果出发对实验数据进行分析。首先，从整体趋势来看，本文方法在两个数据集上的秩相关性表现均处于较高水平，其中在 TVSum 上取得了 0.280 的最高 Spearman 系数，说明所提出的层级语义评分机制能够较为准确地恢复人工标注对帧重要性的全局排序结构。相较于 dppLSTM、VASNet、DSNet 和 iPTNet 等主要依赖视觉表征回归或时序建模的监督方法，本文方法在 $\rho$ 与 $\tau$ 指标上均表现出更明显的优势，表明仅依赖视觉显著性或局部时序依赖难以充分刻画帧级重要性与视频整体主题之间的对应关系，而通过视频级主题摘要-场景级语义评分-帧级语义增强的逐层传递，可以在无需参数训练的前提下提升模型输出与人工排序之间的一致性。

进一步地，与 DMASum、MSVA、CSTA 及 LLMVS 等相关性表现较强的方法相比，本文方法呈现出较为鲜明的特征：在 SumMe 上，本文方法的 $\rho=0.267$、$\tau=0.218$，虽略低于 LLMVS 和 CSTA，但仍显著优于多数传统监督方法；在 TVSum 上，本文方法的 $\rho$ 超过了 DMASum、CSTA 与 LLMVS，仅在 $\tau$ 上略低于 LLMVS 和 DMASum。该结果说明，本文方法在主题结构更清晰、叙事主线更集中的视频中，更擅长把握帧级重要性的整体单调排序趋势，而在更强调局部成对顺序一致性的 Kendall 指标上仍存在进一步优化空间。综合而言，秩相关性实验结果与前述 $F_1$ 分数实验相互印证，共同说明本文方法不仅能够在摘要选择结果上取得较强竞争力，而且能够在帧级评分层面形成与人工判断较为一致的语义排序结构。

\subsubsection{定性分析}
为了从可解释性角度进一步观察模型行为，本文从 SumMe 与 TVSum 数据集中各选取一段视频，对帧级评分曲线与关键帧选择结果进行可视化分析。图~\ref{fig:vis_playing_ball} 与图~\ref{fig:vis_WxtbjNsCQ8A} 分别展示了两段视频的帧级重要性评分曲线及代表性关键帧。每幅图的上方展示模型评分与人工标注取得较高一致性的 6 个时刻所对应的视频帧缩略图；下方由上至下依次绘制场景级语义评分 $S_{\mathrm{LLM}}$（蓝色）、帧--场景相似度 $S_{FS}$（绿色）、帧--视频全局相似度 $S_{FV}$（红色）以及最终融合评分 $F(p_j)$（黑色）四条曲线，各子图中以灰色填充区域标示人工标注的重要性分布，便于直观比较模型评分与参考标注之间的对应关系。

\begin{figure}
  \centering
  \includegraphics[width=\textwidth]{c3/method1_results_visualization_playing_ball.pdf}
  \caption{SumMe 数据集 Playing Ball 视频的帧级评分曲线与代表性关键帧可视化}\label{fig:vis_playing_ball}
\end{figure}

\begin{figure}
  \centering
  \includegraphics[width=\textwidth]{c3/method1_results_visualization_WxtbjNsCQ8A.pdf}
  \caption{TVSum 数据集 WxtbjNsCQ8A 视频的帧级评分曲线与代表性关键帧可视化}\label{fig:vis_WxtbjNsCQ8A}
\end{figure}

\textbf{SumMe Playing Ball 视频分析。}\enspace 如图~\ref{fig:vis_playing_ball} 所示，该视频记录了动物与人之间的日常互动，整体叙事较为松散，属于典型的开放域视频。对此，$S_{\mathrm{LLM}}$ 呈现出明显的分段阶梯变化，对猫与球互动、犬只追逐等核心事件赋予较高场景分数，而对静态空镜与低动态过渡片段给出较低评分，说明所提方法能够较好地区分语义关键内容与冗余内容。在场景内部，$S_{FS}$ 进一步突出动作完整、主体清晰的代表性帧，$S_{FV}$ 则为与视频主旨高度一致但处于中低分场景中的局部关键帧提供补偿。最终融合评分 $F(p_j)$ 的高分区域与人工标注的重要性分布基本一致，表明本文方法在开放域视频中仍能较稳定地定位关键内容并抑制冗余片段。

\textbf{TVSum WxtbjNsCQ8A 视频分析。}\enspace 如图~\ref{fig:vis_WxtbjNsCQ8A} 所示，该视频围绕养蜂人的日常劳作展开，主题明确、叙事主线清晰。与 SumMe 案例相比，$S_{\mathrm{LLM}}$ 在该视频中表现出更清楚的由低到高的层次变化：开场铺垫片段得分较低，进入蜂箱管理、蜂框检查和蜂蜜采集等核心环节后，场景分数明显上升并持续保持高位，较好反映了视频叙事重心的推进过程。同时，$S_{FS}$ 能够在高分场景内筛出更具代表性的操作帧，$S_{FV}$ 则持续强化与养蜂主题高度相关的关键画面。最终融合评分在核心劳动片段上形成连续高响应区间，与人工标注的高分区域较为一致，说明本文方法在主题集中的长视频中具有较好的排序稳定性与关键帧识别能力。

综合来看，两段视频的可视化结果表明，本文方法既能在开放域场景中识别关键事件，也能在主题明确的视频中保持较好的叙事对齐能力。

\xsubsection{消融实验}{Ablation Studies}
为进一步说明本文框架中各组成部分的具体贡献，本节从模块消融分析与参数分析两个层面展开讨论。前者关注不同功能模块及评分策略对整体性能的影响，后者考察关键超参数对模型表现的敏感性。

\subsubsection{模块消融分析}
本文进一步对各功能模块进行消融研究。表\ref{tab:ablation_all} 采用逐次移除单一模块的方式，直观展示层级语义描述、场景语义评分以及帧级分数增强等关键组成部分对完整模型性能的贡献。

\begin{table}
  \centering
  \caption{各子模块消融实验结果}\label{tab:ablation_all}
  \begin{tblr}{
    width=\linewidth,
    colspec={X[3,l]X[c]X[c]X[c]X[c]},
    rowsep=2pt,
    cell{1}{2}={c=2}{c},
    cell{1}{4}={c=2}{c},
    % hline{3}={2-3}{0.5pt},
    % hline{3}={4-5}{0.5pt},
  }
    \toprule
    \textbf{模块} & \textbf{SumMe} & & \textbf{TVSum} & \\
    \midrule
    & Max $F_1$ (\%) & $\Delta$ & Avg $F_1$ (\%) & $\Delta$ \\
    \midrule
    w/o Caption \& Scene Scoring & 45.8 & $-2.3$ & 58.9 & $-2.2$ \\
    w/o $S_{FS}$ & 46.7 & $-1.4$ & 60.6 & $-0.5$ \\
    w/o $S_{FV}$ & 42.2 & $-5.9$ & 60.3 & $-0.8$ \\
    w/o Frame Refinement & 42.2 & $-5.9$ & 52.8 & $-8.3$ \\
    \midrule
    完整方法 & \textbf{48.1} & --- & \textbf{61.1} & --- \\
    \bottomrule
  \end{tblr}
\end{table}

从表\ref{tab:ablation_all} 可以看出，多级语义描述与场景语义评分模块能够为最终摘要提供稳定的高层语义先验。当移除 Caption 与 Scene Scoring 模块后，模型在 SumMe 与 TVSum 上分别下降 2.3 和 2.2 个百分点，说明若缺少由帧级描述、场景级聚合和视频级概括构成的层级语义链路，后续帧级筛选将难以获得可靠的主题参照，摘要质量会出现明显退化。

从帧级增强模块内部两个分支的消融结果可以看出，$S_{FV}$ 的贡献更为关键。移除该分支后，模型在 SumMe 和 TVSum 上分别下降 5.9 和 0.8 个百分点，表明帧与视频全局主题之间的直接语义关联是维持摘要质量的重要因素。相比之下，移除 $S_{FS}$ 后，模型在 SumMe 和 TVSum 上分别下降 1.4 和 0.5 个百分点，说明帧与场景之间的局部语义匹配同样有效，但其贡献幅度整体弱于 $S_{FV}$。这表明在开放域视频中，全局主题约束对关键帧筛选更稳定，而局部场景匹配主要起到细化场景内部排序的辅助作用。

当同时移除 $S_{FS}$ 与 $S_{FV}$、仅保留场景级粗粒度分数作为最终评分时，模型在 TVSum 上进一步下降至 52.8\%，而在 SumMe 上则与移除 $S_{FV}$ 时持平，均为 42.2\%。这一结果说明，若缺乏帧级分辨率的语义细化，模型虽然仍能大致判断哪些场景重要，但难以在高分场景内部完成精确的关键帧甄别。综合来看，层级语义描述与场景评分模块为整套方法提供了必要的主题建模能力，而 $S_{FS}$ 与 $S_{FV}$ 则在帧级层面形成互补，共同支撑最终的摘要质量。

\subsubsection{参数分析}
本文进一步考察平衡系数 $\alpha$ 对模型性能的影响。根据评分公式
\begin{equation}
  F(p_j)=\alpha\,S_{\mathrm{LLM}}(\sigma(j))\,S_{FS}(p_j)+S_{FV}(p_j),
\end{equation}
$\alpha$ 用于调节场景级语义先验与全局主题匹配项之间的相对权重。为分析其敏感性，本文在 $[0.1,1.0]$ 区间内以 0.1 为步长进行扫描，实验结果如图~\ref{fig:alpha_f1_curves} 所示。
\begin{figure}[tbp]
  \centering
  \includegraphics[width=\linewidth]{c3/alpha_f1_curves.pdf}
  \caption{参数 $\alpha$ 对 SumMe 与 TVSum 数据集上 $F_1$ 分数的影响}\label{fig:alpha_f1_curves}
\end{figure}
可以发现，当 $\alpha$ 由 0.1 增加到 0.2 时，SumMe 数据集上的 $F_1$ 分数由 47.50\% 提升至 48.10\%，达到最高值；TVSum 数据集上的性能则在 $\alpha=0.3$ 时达到最高，为 61.05\%。整体来看，两个数据集在 $\alpha\in[0.1,0.3]$ 区间内均保持较高性能，说明适度引入场景级语义先验有助于提升摘要质量。

当 $\alpha$ 继续增大时，两个数据集上的性能均呈下降趋势，其中 TVSum 上的下降更为连续和明显。这主要是因为随着场景先验权重不断提高，评分公式中的乘积项逐渐占据更大比重，从而削弱了全局主题匹配项 $S_{FV}$ 对关键帧筛选的直接约束作用。进一步地，当 $\alpha$ 取值过大时，场景内部帧间差异会被压缩，模型对细粒度重要性排序的刻画能力随之下降。由此可以看出，较小的 $\alpha$ 更有利于在场景先验约束与全局主题匹配之间取得平衡，从而获得更稳定的摘要性能。

\xsection{本章小结}{Chapter Summary}
本章针对开放域视频摘要任务中语义表达不足、对人工标注依赖较强以及跨场景泛化能力有限等问题，提出了一种基于大语言模型的免训练视频摘要方法。该方法以视频内容的层级语义组织为核心，将视频摘要由传统的视觉代表性筛选问题转化为面向全局主题一致性的语义建模问题，从而为开放域场景下的视频摘要生成提供了一种无需参数训练的解决思路。

围绕上述目标，本文构建了由视频场景划分、多级语义描述生成、全局主题约束下的场景语义评分以及场景先验引导的帧级语义分数增强组成的完整框架。该框架首先利用 KTS 将长视频组织为时间连续的场景单元，再借助视觉语言模型和大语言模型逐步形成帧级、场景级与视频级语义表示；在此基础上，通过场景级语义评分建立局部内容与全局主题之间的对应关系，并结合 $S_{FS}$ 与 $S_{FV}$ 对帧级重要性进行细化，从而实现由全局主题到局部关键帧的层级化摘要生成。

实验结果表明，所提方法在 SumMe 和 TVSum 数据集上分别取得了 48.1\% 和 61.1\% 的 $F_1$ 分数，并在 TVSum 上取得了 0.280 的 Spearman 秩相关系数，体现出较好的摘要选择能力与帧级排序一致性。消融实验进一步验证了多级语义描述、场景语义评分以及帧级语义增强等关键模块的必要性，参数分析说明平衡系数 $\alpha$ 取较小值时更有利于兼顾场景先验与全局主题约束。本文方法在不引入任何可学习参数的前提下，有效提升了视频摘要的语义理解能力、结果可解释性与开放域适应性。
